\section{Related Work}
\label{related_work}

\subsection{Traditional Requirement Modeling}
\label{sec:rw-traditional}

Traditional requirements engineering has developed modeling notations, templates, and controlled languages to reduce ambiguity and improve stakeholder communication.
Controlled Natural Language approaches show that structured expression can improve analyzability while preserving readability~\cite{darifControlledNaturalLanguage2026}.
Our work follows this motivation but focuses on automatically constructing structured representations from informal natural-language requirements.

\subsection{Requirement Representation and Requirement Graphs}
\label{sec:rw-representation}

Requirement representation research studies how requirement elements and their relations can be made explicit.
Graph-based representations are useful because they expose entities, dependencies, conditions, and semantic links in a form that can be inspected and analyzed.
This paper treats the requirement graph as the central artifact of requirement structure modeling.

\subsection{LLM-based Requirement Analysis}
\label{sec:rw-llm-requirements}

LLM-based methods have been used to summarize, extract, and organize requirement-related artifacts.
These methods demonstrate the value of making requirements explicit before using them in later engineering stages.
However, fluent generated text does not by itself provide a stable representation for analysis, comparison, or human review.

\subsection{Knowledge-Guided Requirement Modeling}
\label{sec:rw-knowledge}

Knowledge-guided modeling uses domain vocabulary, requirement templates, semantic patterns, and modeling constraints to guide requirement interpretation.
Such knowledge can reduce uncontrolled variation in extracted structures and help align agent outputs with requirements engineering concepts.
This paper uses knowledge as guidance for representation construction rather than as a separate reasoning system.

\subsection{Multi-Agent Collaboration}
\label{sec:rw-multi-agent}

Multi-agent collaboration decomposes complex analysis work into roles with different responsibilities.
For requirement structure modeling, role separation can help distinguish stakeholder-oriented elicitation, semantic decomposition, and graph construction.
This paper investigates a focused three-role workflow instead of a broad multi-agent software engineering platform.
